% GENERATED FILE. Edit canonical lessons, then run npm run generate:books.
% canonical-source-hash: fnv1a64:f6264d5d067c7b28
% canonical-lessons: GE-C454-wanderung, GE-C454-wanderweg, GE-C454-besteigen, GE-C454-absteigen, GE-C454-rasten

\chapter{Eine Wanderung}
\label{ch:ge-eine-wanderung}
\hlchaptermodality{454}

\begin{chapteropening}
\textbf{By the end of this chapter:} \emph{I can talk about a hike: the hike, the trail, climbing up, climbing down and stopping for a rest.}

Complete the last lesson of chapter 454: I can talk about a hike: the hike, the trail, climbing up, climbing down and stopping for a rest.
\end{chapteropening}

\section[die Wanderung]{die Wanderung — the hike, the walk}

\label{lesson:GE-C454-wanderung}



\begin{quote}
Before the new one: say the German for the slope, then the German for the gorge.
\end{quote}

\subsection*{You'll want to know: die Wanderung}
\textbf{die Wanderung} — \textquotedblleft{}the hike, the walk\textquotedblright{}.

\begin{quote}
Die Wanderung zum See dauert drei Stunden.
\end{quote}

\begin{quote}
\emph{The hike to the lake takes three hours.}
\end{quote}

Plural: \emph{die Wanderungen}.

\begin{cousinweb}
\emph{Wanderung} is built on \emph{wandern} 'to hike, to wander' (EtymDB), inherited from Middle High German \emph{wandern} and Old High German \emph{wantarōn}.
\end{cousinweb}

\begin{grammarlens}[title={What you've built}]
One more word for telling, explaining and describing at B1.
\end{grammarlens}

\subsection*{Your turn}
\begin{itemize}
\raggedright
  \item \emph{Say it:} \emph{die Wanderung}
  \item \emph{Say it:} \emph{die Wanderung}, once more
  \item \emph{Say it:} a sentence of your own with \emph{die Wanderung}
  \item \emph{Recall:} say the German for the slope, then the German for the gorge, then say \emph{die Wanderung} again
\end{itemize}

\begin{tcolorbox}[breakable,colback=teal!4,colframe=teal!35!black,title={Before you move on}]
What does die Wanderung mean? (\textquotedblleft{}The hike, the walk\textquotedblright{}.) Say it once more.
\end{tcolorbox}
\section[der Wanderweg]{der Wanderweg — the trail, the footpath}

\label{lesson:GE-C454-wanderweg}



\begin{quote}
Before the new one: say the German for the gorge, then the German for the hike.
\end{quote}

\subsection*{You'll want to know: der Wanderweg}
\textbf{der Wanderweg} — \textquotedblleft{}the trail, the footpath\textquotedblright{}.

\begin{quote}
Der Wanderweg ist mit roten Zeichen markiert.
\end{quote}

\begin{quote}
\emph{The trail is marked with red signs.}
\end{quote}

Plural: \emph{die Wanderwege}.

\begin{cousinweb}
\emph{Wanderweg} is built from \emph{wandern} 'to hike' and \emph{Weg} 'way', inherited from Proto-Germanic \emph{wegaz} (EtymDB).
\end{cousinweb}

\begin{grammarlens}[title={What you've built}]
One more word for telling, explaining and describing at B1.
\end{grammarlens}

\subsection*{Your turn}
\begin{itemize}
\raggedright
  \item \emph{Say it:} \emph{der Wanderweg}
  \item \emph{Say it:} \emph{der Wanderweg}, once more
  \item \emph{Say it:} a sentence of your own with \emph{der Wanderweg}
  \item \emph{Recall:} say the German for the gorge, then the German for the hike, then say \emph{der Wanderweg} again
\end{itemize}

\begin{tcolorbox}[breakable,colback=teal!4,colframe=teal!35!black,title={Before you move on}]
What does der Wanderweg mean? (\textquotedblleft{}The trail, the footpath\textquotedblright{}.) Say it once more.
\end{tcolorbox}
\section[besteigen]{besteigen — to climb (a mountain); to board}

\label{lesson:GE-C454-besteigen}



\begin{quote}
Before the new one: say the German for the hike, then the German for the trail.
\end{quote}

\subsection*{You'll want to know: besteigen}
\textbf{besteigen} — \textquotedblleft{}to climb (a mountain); to board\textquotedblright{}.

\begin{quote}
Im Sommer besteigen sie einen hohen Berg.
\end{quote}

\begin{quote}
\emph{In summer they climb a high mountain.}
\end{quote}

Inseparable, with an object: \emph{sie hat den Gipfel bestiegen}.

\begin{cousinweb}
\emph{besteigen} is built from \emph{be-} and \emph{steigen} 'to climb', inherited from Old High German \emph{stīgan} (EtymDB).
\end{cousinweb}

\begin{grammarlens}[title={What you've built}]
One more word for telling, explaining and describing at B1.
\end{grammarlens}

\subsection*{Your turn}
\begin{itemize}
\raggedright
  \item \emph{Say it:} \emph{besteigen}
  \item \emph{Say it:} \emph{besteigen}, once more
  \item \emph{Say it:} a sentence of your own with \emph{besteigen}
  \item \emph{Recall:} say the German for the hike, then the German for the trail, then say \emph{besteigen} again
\end{itemize}

\begin{tcolorbox}[breakable,colback=teal!4,colframe=teal!35!black,title={Before you move on}]
What does besteigen mean? (\textquotedblleft{}To climb (a mountain); to board\textquotedblright{}.) Say it once more.
\end{tcolorbox}
\section[absteigen]{absteigen — to climb down; to get off (a bike); to go down (a league)}

\label{lesson:GE-C454-absteigen}



\begin{quote}
Before the new one: say the German for the trail, then the German for to climb.
\end{quote}

\subsection*{You'll want to know: absteigen}
\textbf{absteigen} — \textquotedblleft{}to climb down; to get off (a bike); to go down (a league)\textquotedblright{}.

\begin{quote}
Am Nachmittag steigen wir langsam ins Tal ab.
\end{quote}

\begin{quote}
\emph{In the afternoon we slowly climb down into the valley.}
\end{quote}

Separable; Perfekt with \emph{sein}: \emph{wir sind abgestiegen}.

\begin{cousinweb}
\emph{absteigen} is inherited from Middle High German \emph{abestīgen} (EtymDB), built from \emph{ab} 'down' and \emph{steigen} 'to climb'; Dutch \emph{afstijgen} is a cousin.
\end{cousinweb}

\begin{grammarlens}[title={What you've built}]
One more word for telling, explaining and describing at B1.
\end{grammarlens}

\subsection*{Your turn}
\begin{itemize}
\raggedright
  \item \emph{Say it:} \emph{absteigen}
  \item \emph{Say it:} \emph{absteigen}, once more
  \item \emph{Say it:} a sentence of your own with \emph{absteigen}
  \item \emph{Recall:} say the German for the trail, then the German for to climb, then say \emph{absteigen} again
\end{itemize}

\begin{tcolorbox}[breakable,colback=teal!4,colframe=teal!35!black,title={Before you move on}]
What does absteigen mean? (\textquotedblleft{}To climb down; to get off (a bike); to go down (a league)\textquotedblright{}.) Say it once more.
\end{tcolorbox}
\section[rasten]{rasten — to rest, to stop for a rest (on a walk)}

\label{lesson:GE-C454-rasten}



\begin{quote}
Before the new one: say the German for to climb, then the German for to climb down.
\end{quote}

\subsection*{You'll want to know: rasten}
\textbf{rasten} — \textquotedblleft{}to rest, to stop for a rest (on a walk)\textquotedblright{}.

\begin{quote}
Auf halbem Weg rasten wir an einer Hütte.
\end{quote}

\begin{quote}
\emph{Halfway we stop for a rest at a hut.}
\end{quote}

\emph{Die Rast} — the rest. \emph{Rasten} sounds a little old-fashioned; people also say \emph{eine Pause machen}.

\begin{cousinweb}
\emph{rasten} is built on \emph{Rast} 'rest, a break' (EtymDB lists the verb 'to rest'), an old Germanic word; English \emph{rest} is usually counted as a cousin.
\end{cousinweb}

\begin{grammarlens}[title={What you've built}]
One more word for telling, explaining and describing at B1.
\end{grammarlens}

\subsection*{Your turn}
\begin{itemize}
\raggedright
  \item \emph{Say it:} \emph{rasten}
  \item \emph{Say it:} \emph{rasten}, once more
  \item \emph{Say it:} a sentence of your own with \emph{rasten}
  \item \emph{Recall:} say the German for to climb, then the German for to climb down, then say \emph{rasten} again
\end{itemize}

\begin{tcolorbox}[breakable,colback=teal!4,colframe=teal!35!black,title={Before you move on}]
What does rasten mean? (\textquotedblleft{}To rest, to stop for a rest (on a walk)\textquotedblright{}.) Say it once more.
\end{tcolorbox}
